% Neurolib starter preamble
% Curated for graph algorithms / network science / linear algebra notes.
% Bumping a substantive change here invalidates every cached PDF (see PREAMBLE_VERSION).

\usepackage[utf8]{inputenc}
\usepackage[T1]{fontenc}
\usepackage{lmodern}
\usepackage{microtype}

% Math
\usepackage{amsmath}
\usepackage{amssymb}
\usepackage{amsthm}
\usepackage{mathtools}

% Theorem environments
\theoremstyle{definition}
\newtheorem{definition}{Definition}[section]
\newtheorem{example}[definition]{Example}

\theoremstyle{plain}
\newtheorem{theorem}[definition]{Theorem}
\newtheorem{lemma}[definition]{Lemma}
\newtheorem{proposition}[definition]{Proposition}
\newtheorem{corollary}[definition]{Corollary}

\theoremstyle{remark}
\newtheorem*{remark}{Remark}
\newtheorem*{note}{Note}

% TikZ / PGF for diagrams
\usepackage{tikz}
\usepackage{pgfplots}
\usepackage{tikz-cd}
\pgfplotsset{compat=1.18}
% Note: graphdrawing libraries (\usegdlibrary{layered, force, trees}) require
% LuaTeX. Tectonic uses XeTeX, so they're not loaded here. The `graphs` library
% gives you declarative graph syntax; for auto-layout, position nodes manually
% with `positioning`/`calc` or switch to a TeX Live + lualatex toolchain.
\usetikzlibrary{
  graphs,
  arrows.meta,
  positioning,
  calc,
  decorations.pathreplacing,
  shapes,
  shapes.geometric,
  quotes,
  fit,
  matrix
}

% References / hyperlinks (load near the end)
\usepackage[hidelinks]{hyperref}
\usepackage{cleveref}

% Page geometry — comfortable single-column reading
\usepackage[a4paper, margin=1in]{geometry}
\setlength{\parskip}{0.6em}
\setlength{\parindent}{0pt}

% --- Neurolib brand --------------------------------------------------------
\usepackage{xcolor}

\definecolor{neuroink}{HTML}{1F2933}
\definecolor{neurocoral}{HTML}{FF7A59}
\definecolor{neuromuted}{HTML}{5C6470}
\definecolor{neurosky}{HTML}{4DB8E8}
\definecolor{neuromint}{HTML}{3CCB7F}
\definecolor{neurosunshine}{HTML}{FFD93D}
\definecolor{neurocream}{HTML}{FFFDF7}

% Logo: an open book viewed from above with a small knowledge graph on each
% page. Recreates public/logo.svg in TikZ so no image asset is needed.
% Optional argument scales it (default 1cm unit).
\newcommand{\neurolibLogo}[1][1]{%
  \begin{tikzpicture}[
    x=#1 cm, y=#1 cm,
    line cap=round, line join=round
  ]
    % Pages
    \filldraw[fill=neurocream, draw=neuroink, line width=0.6pt]
      (-1.50, -0.85) -- (-0.05, -1.00) -- (-0.05, 1.00) -- (-1.45, 0.85) -- cycle;
    \filldraw[fill=neurocream, draw=neuroink, line width=0.6pt]
      ( 1.50, -0.85) -- ( 0.05, -1.00) -- ( 0.05, 1.00) -- ( 1.45, 0.85) -- cycle;
    % Spine
    \draw[neuroink, line width=0.6pt] (0, -1.00) -- (0, 1.00);

    % Page-internal edges
    \draw[neuromuted, line width=0.3pt]
      (-1.10, 0.50) -- (-0.70, 0) -- (-1.05, -0.50) -- cycle;
    \draw[neuromuted, line width=0.3pt]
      ( 1.10, 0.50) -- ( 0.70, 0) -- ( 1.05, -0.50) -- cycle;

    % Spine-crossing dashed edges
    \draw[neurocoral, dash pattern=on 0.4pt off 0.8pt, line width=0.3pt]
      (-1.10, 0.50) -- (0, 0.70);
    \draw[neurocoral, dash pattern=on 0.4pt off 0.8pt, line width=0.3pt]
      (-0.70, 0)    -- (0, 0);
    \draw[neurocoral, dash pattern=on 0.4pt off 0.8pt, line width=0.3pt]
      (-1.05, -0.50) -- (0, -0.70);
    \draw[neurocoral, dash pattern=on 0.4pt off 0.8pt, line width=0.3pt]
      ( 1.10, 0.50) -- (0, 0.70);
    \draw[neurocoral, dash pattern=on 0.4pt off 0.8pt, line width=0.3pt]
      ( 0.70, 0)    -- (0, 0);
    \draw[neurocoral, dash pattern=on 0.4pt off 0.8pt, line width=0.3pt]
      ( 1.05, -0.50) -- (0, -0.70);

    % All nodes: coral
    \filldraw[fill=neurocoral, draw=neuroink, line width=0.25pt]
      (0,  0.70) circle (0.085);
    \filldraw[fill=neurocoral, draw=neuroink, line width=0.25pt]
      (0,  0)    circle (0.075);
    \filldraw[fill=neurocoral, draw=neuroink, line width=0.25pt]
      (0, -0.70) circle (0.085);
    \filldraw[fill=neurocoral, draw=neuroink, line width=0.25pt]
      (-1.10,  0.50) circle (0.080);
    \filldraw[fill=neurocoral, draw=neuroink, line width=0.25pt]
      (-0.70,  0)    circle (0.080);
    \filldraw[fill=neurocoral, draw=neuroink, line width=0.25pt]
      (-1.05, -0.50) circle (0.080);
    \filldraw[fill=neurocoral, draw=neuroink, line width=0.25pt]
      ( 1.10,  0.50) circle (0.080);
    \filldraw[fill=neurocoral, draw=neuroink, line width=0.25pt]
      ( 0.70,  0)    circle (0.080);
    \filldraw[fill=neurocoral, draw=neuroink, line width=0.25pt]
      ( 1.05, -0.50) circle (0.080);
  \end{tikzpicture}%
}

% Cover page used by every compile (single card + composition).
% Args: 1 = title, 2 = subtitle (subject path or composition kind).
\newcommand{\neurolibCover}[2]{%
  \begin{titlepage}
    \centering
    \vspace*{2cm}%
    \neurolibLogo[1.6]%
    \par\vspace{0.6em}%
    {\color{neurocoral}\textsc{\fontsize{14pt}{16pt}\selectfont Neurolib}}%
    \par\vspace{3.5cm}%
    {\fontsize{32pt}{38pt}\selectfont\bfseries\color{neuroink} #1\par}%
    \vspace{1.2cm}%
    {\color{neuromuted}\fontsize{15pt}{19pt}\selectfont #2\par}%
    \vfill
    {\color{neuromuted}\fontsize{11pt}{13pt}\selectfont \today\par}%
    \vspace{2cm}%
  \end{titlepage}%
}
